\title{QLoRA: Efficient Finetuning of Quantized LLMs}

\begin{abstract}
We introduce \textsc{QLoRA}, a finetuning method designed to optimize memory efficiency, allowing the finetuning of a 65B parameter model using a single 48GB GPU without compromising the performance typically associated with full 16-bit finetuning. \textsc{QLoRA} enables gradient backpropagation through a static, 4-bit quantized pretrained language model directly into Low Rank Adapters~(LoRA). Our premier model suite, termed \textbf{Guanaco}, surpasses all prior openly available models on the Vicuna benchmark, achieving 99.3\% of ChatGPT’s performance after just 24 hours of training on a single GPU. Key innovations introduced by \textsc{QLoRA} for reducing memory consumption while maintaining performance include: (a) 4-bit NormalFloat (NF4), a new datatype optimized from an information theory perspective for normally distributed weights; (b) Double Quantization, which further minimizes memory usage by quantizing the quantization parameters themselves; and (c) Paged Optimizers, which mitigate memory usage spikes during optimization. By applying \textsc{QLoRA}, we finetune over 1,000 models, enabling comprehensive analysis of instruction-following and chatbot capabilities across 8 instruction datasets, various architectures (LLaMA, T5), and model sizes that would be impractical with traditional finetuning, such as 33B and 65B parameter models. Our findings reveal that applying QLoRA finetuning on compact, high-quality datasets consistently yields state-of-the-art performance, even with smaller architectures compared to the previous best-performing models. We offer in-depth evaluation of chatbot efficacy via both human and GPT-4 assessment, demonstrating that GPT-4 can serve as a cost-effective and reliable alternative to human evaluation. Additionally, we show that existing chatbot benchmarks are insufficient for reliably measuring chatbot competence. A targeted analysis highlights specific instances where \textbf{Guanaco} falls short relative to ChatGPT. We make all models and source code publicly available, including CUDA kernels supporting 4-bit training.\footnote{\url{https://github.com/artidoro/qlora} and \url{https://github.com/TimDettmers/bitsandbytes}}
\end{abstract}

\section{Introduction}

Finetuning large language models (LLMs) has proven to be a powerful approach for enhancing their capabilities \citep{min2021metaicl, wei2021finetuned, ouyang2022training, wang2022super, wang2022self, liu2022few}, as well as for encouraging or discouraging particular behaviors \citep{ouyang2022training,askell2021general,bai2022training}. Nonetheless, the computational and memory costs associated with finetuning extremely large models are prohibitive; for instance, standard 16-bit finetuning of the LLaMA model with 65 billion parameters~\citep{touvron2023llama} necessitates over 780 GB of GPU memory. Although emerging quantization strategies can mitigate the memory usage of LLMs during inference \citep{dettmers2022llm, dettmers2022case, frantar2022gptq, xiao2022smoothquant}, these methods typically fail when applied to the training process \citep{wortsman2023stable}.

In this work, we establish for the first time that it is feasible to finetune LLMs at 4-bit quantization with no loss in performance. Our proposed method, \textsc{QLoRA}, introduces a novel high-precision quantization approach that compresses a pretrained model down to 4 bits and incorporates a compact set of trainable Low-rank Adapter weights \citep{hu2021lora} 

which are updated by propagating gradients through the quantized model parameters.

By employing \textsc{QLoRA}, the average memory demand for finetuning a 65B parameter model drops dramatically from over 780 GB to less than 48 GB, all while maintaining the runtime and predictive accuracy relative to the full 16-bit finetuning baseline. This breakthrough substantially lowers the barrier for LLM finetuning, making it possible to finetune the largest public models on a single GPU. Leveraging \textsc{QLoRA}, we introduce the \textbf{Guanaco} suite of models, with the second largest achieving 97.8\% of ChatGPT's performance on the Vicuna~\citep{vicuna2023} benchmark, and training can be completed in under 12 hours on an individual consumer GPU; with a single professional GPU and a 24-hour training run, our largest model attains 99.3\% of ChatGPT's performance, effectively bridging the remaining gap on the Vicuna benchmark. Upon deployment, our smallest \textbf{Guanaco} model (7B parameters) needs only 5 GB of memory and surpasses a 26 GB Alpaca model by more than 20 percentage points on the Vicuna benchmark (see Table~\ref{tbl:vicuna_eval}).

\textsc{QLoRA} incorporates several novel techniques aimed at minimizing memory consumption while maintaining model efficacy: (1) \textbf{4-bit NormalFloat}, an optimal quantization format—rooted in information theory—for handling normally distributed data, which empirically surpasses both 4-bit Integers and 4-bit Floats. (2) \textbf{Double Quantization}, which involves further quantizing the quantization constants themselves, achieving an average reduction of roughly 0.37 bits per parameter (equivalent to about 3 GB saved for a 65B parameter model). (3) \textbf{Paged Optimizers}, which utilize NVIDIA's unified memory to mitigate the memory spikes traditionally seen with gradient checkpointing when managing mini-batches with extensive sequence lengths. By integrating these advancements, we present a refined LoRA method that features adapter modules at each network layer, nearly eliminating the accuracy compromises previously observed.

The memory efficiency provided by \textsc{QLoRA} allows for a comprehensive exploration of instruction finetuning and chatbot efficacy on model scales otherwise infeasible with standard finetuning due to memory limitations. Consequently, we train in excess of 1,000 models across a variety of instruction tuning datasets, architectures, and parameter counts ranging from 80M to 65B. Beyond demonstrating that \textsc{QLoRA} attains 16-bit precision performance (\S\ref{sec:comp_to_std}) and producing a cutting-edge chatbot, \textbf{Guanaco} (\S\ref{sec:pushing}), we further examine emerging patterns in the trained models. Notably, our findings indicate that the quality of training data has a substantially greater impact than sheer dataset magnitude—for instance, a compact dataset of 9k samples (OASST1) surpassed a much larger 450k sample dataset (FLAN v2, subsampled) in chatbot tasks, despite both being curated for instruction-following generalization. Furthermore, we observe that excelling at the Massive Multitask Language Understanding (MMLU) benchmark does not necessarily correspond with high performance on the Vicuna chatbot benchmark, underscoring that aligning dataset characteristics with the target task takes precedence over merely expanding dataset size.

Additionally, we undertake a thorough assessment of chatbot effectiveness employing both human evaluators and GPT-4 as judges. Our evaluation utilizes a tournament-style framework in which different models are pitted against one another, responding to prompts in direct matchups. For each match, either GPT-4 or human raters determine which response is superior. The results across all matchups are synthesized using Elo ratings~\citep{elo1967proposed, elo1978rating}, enabling an ordinal ranking of the models’ capabilities. Our analysis demonstrates that rankings derived from GPT-4 and those from human evaluators are largely concordant, although notable disagreements do occur in specific cases. Consequently, we emphasize that while model-based evaluation offers a cost-efficient surrogate for human annotation, it comes with inherent limitations and sources of variance.

To further contextualize our quantitative benchmarks, we supplement our evaluation of \textbf{Guanaco} models with qualitative insights. This deeper analysis uncovers strengths and weaknesses that escape measurement by purely quantitative means.

For the benefit of the community, we are releasing all generated model outputs along with corresponding human and GPT-4 assessments. We are also making our full codebase and CUDA implementations publicly available, with integration into the Hugging Face transformers framework \citep{wolf2019huggingface} to enhance usability. Furthermore, we release adapter weights for models with 7/13/33/65B parameters, each trained on 8 distinct instruction-tuning datasets, culminating in 32 unique fine-tuned, open-sourced model variants.

\section{Background}
\paragraph{Block-wise k-bit Quantization} The aim of quantization is to map data from a high-precision format to a more compact, lower-precision representation. Typically, this involves converting variables from a higher bit-width (such as 32-bit floats) to a lower one (like 8-bit integers). To exploit the available numerical range of the target low-precision type, the values are usually normalized based on the absolute maximum value from the source data—commonly represented as a tensor—prior to conversion. For instance, to quantize a tensor of 32-bit floating points (FP32) into an 8-bit integer tensor (Int8) within $[-127, 127]$, one can use:

\begin{equation}
    \mathbf{X}^{ \text{Int8}} = \text{round}\left( \frac{127}{{\text{absmax}}(\mathbf{X}^{{\text{FP32}}})}\mathbf{X}^{{\text{FP32}}} \right) = \text{round}(c^{\text{FP32}}\cdot \mathbf{X}^{{\text{FP32}}}),
\end{equation}

where $c$ denotes the {\em quantization constant} or {\em quantization scale}. The reverse transformation—dequantization—is described as follows:

\begin{equation}
    \text{dequant}(c^{\text{FP32}}, \mathbf{X}^{\text{Int8}}) = \frac{\mathbf{X}^{\text{Int8}}}{c^{\text{FP32}}} = \mathbf{X}^{\text{FP32}}
\end{equation}

A major drawback of this method arises when an outlier—an input tensor element with unusually high magnitude—appears. In such scenarios, the distribution of quantization bins (the mapping to specific bit patterns) becomes inefficient, as many bins may remain unused or sparsely populated. To mitigate the impact of these outliers, practitioners frequently segment the input tensor into several blocks, each undergoing independent quantization and assigned its unique quantization constant $c$. Formally, the input tensor $\mathbf{X}\in \mathbb{R}^{b\times h}$ is reshaped into $n$ contiguous blocks of length $B$ by linearizing the tensor and partitioning it into ${n = ({b\times h} )/{B} }$ segments. Each of these segments is quantized separately according to Equation~1, yielding a quantized representation and a collection of $n$ quantization constants $c_i$.

\paragraph{Low-rank Adapters} The Low-rank Adapter (LoRA) approach to finetuning \citep{hu2021lora} addresses memory constraints by introducing a compact set of trainable adapter parameters, while freezing the rest of the model’s weights. During training, gradients produced by stochastic gradient descent flow through the pretrained (and unaltered) parameters into the adapter, which is subsequently updated to minimize the loss. LoRA enhances a standard linear transformation by injecting an extra, factorized matrix. For a linear transformation given as ${\mathbf{X} \mathbf{W} = \mathbf{Y}}$ with $\mathbf{X}\in  \mathbb{R}^{b\times h}$ and $\mathbf{W}\in  \mathbb{R}^{h\times o}$, the output under LoRA becomes:

\begin{equation}
    \mathbf{Y} = \mathbf{X}\mathbf{W} + s\mathbf{X}\mathbf{L}_1\mathbf{L}_2,
\end{equation}

where $\mathbf{L}_1\in  \mathbb{R}^{h\times r}$ and $\mathbf{L}_2\in  \mathbb{R}^{r\times o}$ are learnable parameters, and $s$ denotes a scaling coefficient.

\paragraph{Memory Requirement of Parameter-Efficient Finetuning} The consideration of LoRA’s memory consumption during training, both in terms of the quantity and dimensionality of adapters employed, is a crucial issue. Due to LoRA's inherently small memory demands, it becomes feasible to introduce additional adapters to enhance model effectiveness without notably increasing total memory usage. Although LoRA is developed as a Parameter Efficient Finetuning (PEFT) approach, the main contributor to memory usage during LLM finetuning is typically the storage of activation gradients rather than the LoRA parameters themselves. For instance, training a 7B LLaMA model on FLAN v2 with a batch size of 1 and applying LoRA with weight allocations corresponding to the usual 0.2\% of original model parameters~\citep{hu2021lora,liu2022few}, results in LoRA input gradients occupying 567 MB of memory, whereas the parameters only require 26 MB. By incorporating gradient checkpointing~\citep{chen2016training}, the average memory needed for input gradients per sequence can be lowered to 18 MB, which remains higher than the memory used by all LoRA weights together. For comparison, the 4-bit quantized base model alone utilizes 5,048 MB. These figures underscore that, while gradient checkpointing is vital for managing memory, substantial reductions in the size of LoRA parameters yield only marginal improvements in overall memory efficiency. Consequently, the number of adapters can be increased without meaningfully expanding memory demands during training (see Appendix~\ref{app:memory} for further details). As will be elaborated subsequently, this property is essential to achieving performance levels comparable to models trained with full 16-bit precision.

\section{\textsc{QLoRA} Finetuning}

\textsc{QLoRA} facilitates accurate finetuning with 4-bit precision through two key innovations introduced in this work—4-bit NormalFloat (NF4) quantization and Double Quantization. Furthermore, to mitigate the risk of out-of-memory errors resulting from memory surges during gradient checkpointing, which has historically hindered single-machine finetuning of large models, we incorporate Paged Optimizers.

The approach utilizes a single low-precision storage data type—typically 4-bit—and a separate computation data type, most often BFloat16. Operationally, every time a \textsc{QLoRA} weight tensor is accessed, it is first dequantized into BFloat16, followed by execution of matrix multiplications at the 16-bit level.

We proceed by elaborating on each component constituting \textsc{QLoRA}, and then provide a formal definition.

\paragraph{4-bit NormalFloat Quantization} The NormalFloat (NF) data type extends the concept of Quantile Quantization \citep{dettmers2022optimizers}, which achieves optimal information-theoretical properties by assigning an equal count of tensor values to every quantization bin. This is accomplished by utilizing the empirical cumulative distribution function to estimate the quantiles within the tensor.

However, a significant drawback of quantile quantization lies in the computational intensity of quantile estimation. Consequently, efficient quantile approximation schemes, such as SRAM quantiles~\citep{dettmers2022optimizers}, are adopted in practice. These approximations, while expedient, can lead to significant quantization errors for outlier values, which tend to be particularly consequential.

Such overheads from detailed quantile calculations and inaccuracies introduced by approximation can be alleviated when the input tensors are distributed according to a distribution determined up to a quantization constant. In these instances, since tensors share identical quantiles, precise quantile estimation is computationally tractable.

Given that the pretrained neural network weights typically follow a normal distribution centered at zero with a standard deviation $\sigma$ (refer to Appendix~\ref{app:normality}), we can standardize all weight distributions by adjusting $\sigma$ so that the resulting distribution aligns precisely with the available data type range. For this study, we designate the target range of our data type as $[-1, 1]$. Thus, it is necessary to map both the data type quantiles and the neural network weights into this predefined interval through normalization.

The data type that is theoretically optimal from an information-theoretic standpoint for handling zero-mean normal distributions with varying standard deviations $\sigma$ constrained to $[-1, 1]$ can be obtained using the following procedure: (1) compute the $2^k+1$ quantile boundaries for the standard normal distribution $N(0, 1)$, thus deriving a $k$-bit quantile quantization scheme suitable for normal distributions, (2) rescale the data type values so they are distributed within $[-1, 1]$, (3) transform a given weight tensor by normalizing it into the $[-1, 1]$ interval via absolute maximum rescaling.

When both the weights and data type values occupy the same normalized range, quantization proceeds as standard. The normalization in step (3) corresponds to scaling the weight tensor's standard deviation to that of the $k$-bit quantized data type. The data type values $q_i$ for $2^k$ quantization bins can be formally determined by:

\begin{equation}
q_i  = \frac{1}{2}\left( Q_X\left(\frac{i}{2^k+1}\right) + Q_X\left(\frac{i+1}{2^k+1}\right)\right),
\end{equation}

where $Q_X(\cdot)$ denotes the quantile function for the standard normal distribution $N(0,1)$. One limitation of a symmetric k-bit quantization scheme is its inability to precisely encode the value zero, which is particularly problematic when zero-valued elements, such as those used in padding, must be represented without error. To guarantee that zero has a discrete representation and to fully utilize the $2^k$ distinct values of a k-bit format, we develop an asymmetric quantization method. This is accomplished by independently computing the quantiles $q_i$ for the negative domain (using $2^{k-1}$ bins) and for the positive domain (with $2^{k-1}+1$ bins). These quantile sets are then combined, with the duplicated zero entry eliminated to preserve a unique representation of zero. We refer to this data format as \emph{k-bit NormalFloat} (NFk), as each quantization interval is expected to contain an equal probability mass under a zero-mean normal distribution, rendering it information-theoretically optimal for such data. The complete specification of this datatype is detailed in Appendix~\ref{app:nf4}.

\paragraph{Double Quantization{}} 
We propose a technique termed \emph{Double Quantization} (DQ), which entails a secondary quantization applied to the quantization coefficients themselves, offering further reductions in memory consumption. While small block sizes are necessary to achieve accurate 4-bit quantization~\citep{dettmers2022case}, they incur non-negligible memory overhead. As an illustration, with 32-bit quantization coefficients and a block size of 64 in the quantization of $\mathbf{W}$, the storage required for these coefficients averages $32/64 = 0.5$ bits per parameter. Double Quantization serves to minimize the additional memory needed for storing these quantization coefficients.

In practical terms, Double Quantization involves taking the quantization constants $c_2^{\text{FP32}}$ produced in the first quantization step and subjecting them to a second quantization process. This secondary quantization yields both quantized versions of the quantization constants $c_2^{\text{FP8}}$ and an additional set of quantization coefficients $c_1^{\text{FP32}}$ for the second stage. For this secondary quantization, we employ 8-bit floating point numbers and a block size of 256, as empirical evidence—including findings from \citet{dettmers2022case}—indicates that such 8-bit quantization does not degrade performance. To facilitate symmetric quantization, the mean of $c_2$ is subtracted prior to quantization, centering the coefficients around zero, since $c_2^{\text{FP32}}$ are inherently positive. With a block size of 64, this approach brings down the memory cost per parameter from $32/64=0.5$ bits to $8/64 + 32/(64\cdot256) = 0.127$ bits, achieving a per-parameter savings of 0.373 bits.

\paragraph{Paged Optimizers} leverage the NVIDIA unified memory~\footnote{\tiny \url{https://docs.nvidia.com/cuda/cuda-c-programming-guide}} capability, which enables automatic transfers of memory pages between the CPU and GPU. This mechanism allows for uninterrupted GPU operations even when the GPU exhausts its available memory. The unified memory system acts analogously to standard CPU memory paging between main memory and disk. In our approach, we utilize this mechanism by allocating the optimizer state memory as paged memory. As a result, when GPU memory is insufficient, the optimizer states are transparently swapped to CPU RAM and subsequently returned to GPU memory on demand during optimizer update computations.

\paragraph{\textsc{QLoRA}.} Based on the outlined system components, we formalize \textsc{QLoRA} for a single linear transformation in the quantized underlying model with an individual LoRA adapter as:

\begin{equation}
    \mathbf{Y}^{\text{BF16}} =  \mathbf{X}^{\text{BF16}} \text{doubleDequant}(c_1^{\text{FP32}}, c_2^{\text{k-bit}}, \mathbf{W}^{\text{NF4}}) + \mathbf{X}^{\text{BF16}}\mathbf{L}^{\text{BF16}}_1\mathbf{L}^{\text{BF16}}_2,
\end{equation}

The function doubleDequant$(\cdot)$ is expressed as:

\begin{equation}
    \text{doubleDequant}(c_1^{\text{FP32}}, c_2^{\text{k-bit}}, \mathbf{W}^{\text{k-bit}}) = \text{dequant}(\text{dequant}(c_1^{\text{FP32}}, c_2^{\text{k-bit}}), \mathbf{W}^{\text{4bit}}) = \mathbf{W}^{\text{BF16}},
\end{equation}

Here, $\mathbf{W}$ is quantized with NF4, and $c_2$ is represented in FP8 format.

To improve quantization accuracy, a blocksize of 64 is chosen for $\mathbf{W}$, while a larger blocksize of 256 is used for $c_2$ to reduce memory footprint.

During training, only gradients of the error with respect to the adapter weights, $\frac{\partial E}{\partial \mathbf{L}_i}$, are required; gradients with respect to the 4-bit quantized weights, $\frac{\partial E}{\partial \mathbf{W}}$, are not needed. Nevertheless, computing $\frac{\partial E}{\partial \mathbf{L}_i}$ requires evaluating $\frac{\partial \mathbf{X}}{\partial \mathbf{W}}$, which is done via equation~(5) by dequantizing the stored $\mathbf{W}^{\text{NF4}}$ into the computation data type $\mathbf{W}^{\text{BF16}}$, so that the derivative $\frac{\partial \mathbf{X}}{\partial \mathbf{W}}$ can be obtained in BFloat16 precision.

In brief, \textsc{QLoRA} operates using a single storage data type, typically 4-bit NormalFloat, alongside a distinct computation data type, specifically 16-bit BrainFloat. During both the forward and backward propagation steps, the parameters stored as the lower-precision data type are dequantized to the higher-precision computation data type; however, gradient computations are performed exclusively for the LoRA-specific parameters, which are maintained in 16-bit BrainFloat.

\section{QLoRA vs. Standard Finetuning}
\label{sec:comp_to_std}

Having outlined the mechanisms of QLoRA and demonstrated its potential for substantial memory savings during model finetuning, we turn to the key issue of whether QLoRA can deliver comparable results to conventional full-model finetuning. Moreover, we aim to dissect QLoRA's design choices, particularly by examining the impact of NormalFloat4 compared to standard Float4 quantization. The following subsections elaborate on experiments designed to address these aspects.

\paragraph{Experimental setup.} Our investigation spans three model categories—encoder, encoder-decoder, and decoder-only architectures—where we benchmark QLoRA against 16-bit adapter-based finetuning and full finetuning for models as large as 3B parameters. Evaluation tasks comprise GLUE \citep{wang2018glue} on RoBERTa-large \citep{liu2019roberta}, Super-NaturalInstructions (TKInstruct) \citep{wang2022super} utilizing T5 \citep{t5}, and 5-shot MMLU \citep{hendrycksmeasuring} following LLaMA finetuning with Flan v2 \citep{longpre2023flan} and Alpaca \citep{alpaca}. To further analyze the benefits of NF4 relative to other 4-bit quantization schemes, we replicate the protocol of \citet{dettmers2022case}—evaluating post-quantization zero-shot accuracy and perplexity across a variety of models (OPT~\citep{zhang2022opt}, LLaMA~\citep{touvron2023llama}, BLOOM~\citep{scao2022bloom}, and Pythia~\citep{biderman2023pythia}) spanning model sizes from 125m up to 13B.

Details tailored to each experimental setting are presented in the respective results sections to enhance clarity, with comprehensive information available in Appendix~\ref{app:exp-setup}.

Paged optimizers play an essential role in enabling \textsc{QLoRA} training for 33B/65B models on GPUs with 24GB or 48GB of memory. Nevertheless, we do not provide direct empirical measurements for Paged Optimizers, as paging events typically only occur when minibatches have especially long sequence lengths—a rare scenario. Notably, in our runtime analysis of paged optimizers with 65B parameter models on 48GB GPUs, using a batch size of 16, we observe that training speeds match those achieved by conventional optimizers. We encourage future work to systematically evaluate scenarios where paging may lead to notable slow-downs.

\paragraph{Default LoRA hyperparameters do not match 16-bit performance} 

Application of the common LoRA approach—targeting the query and value attention projection matrices \citep{hu2021lora}—does not suffice to replicate the performance of full-model finetuning when scaling up to large models. Our findings, illustrated in Figure~\ref{fig:hyper} for LLaMA 7B finetuned on Alpaca, indicate that the critical determinant among LoRA hyperparameters is the number of adapter modules incorporated; applying LoRA to all linear layers in the transformer blocks proves essential for achieving parity with full finetuning results. Other hyperparameters, such as the projection dimension $r$, exhibit minimal effect on performance (as further detailed in Appendix \ref{app:exp-setup}).

In a similar vein, our analysis reveals that the standard hyperparameter configurations employed for fully finetuned baseline models tend to be suboptimal. To establish more competitive baselines, we perform a systematic search across a range of learning rates, from $1\mathrm{e}{-6}$ up to $5\mathrm{e}{-5}$, and batch sizes varying from 8 to 128. The resulting performance outcomes for 7B LLaMA finetuned on Alpaca are presented in Figure~\ref{fig:hyper}.

\paragraph{4-bit NormalFloat Outperforms 4-bit Floating Point}

Although the 4-bit NormalFloat (NF4) format is proven to be information-theoretically optimal, it remains an open question whether these theoretical gains manifest in practical model performance. Adopting the evaluation framework from \citet{dettmers2022case}, we assess quantized large language models—including OPT \citep{zhang2022opt}, BLOOM \cite{scao2022bloom}, Pythia \cite{biderman2023pythia}, and LLaMA—ranging in scale from 125M to 65B parameters. These models, using various quantization types, are evaluated on both language modeling tasks and a collection of zero-shot benchmarks. As depicted in Figure~\ref{fig:nf4} and detailed in Table~\ref{tbl:nf4_ppl}, NF4 not only significantly surpasses FP4 and Int4 in performance but also demonstrates that employing double quantization can further minimize memory usage without incurring a performance penalty.

\paragraph{k-bit \textsc{QLoRA} Equates to 16-bit Full and 16-bit LoRA Finetuning}

Past research has demonstrated that 4-bit quantization can support \emph{inference}, though it typically comes at a cost to performance when compared to 16-bit precision \citep{dettmers2022case,frantar2022gptq}. This prompts a key inquiry: can finetuning adapters at 4-bit precision recover this lost performance? To address this, we conduct experiments under two settings.

The initial setup entails a head-to-head comparison with full 16-bit finetuning, focusing on RoBERTA and T5 models spanning from 125M to 3B parameters. The evaluation is performed on both the GLUE and Super-NaturalInstructions datasets, with findings reported in Table~\ref{tab:nf4vsbf16}. Across both benchmarks, we find that 16-bit, 8-bit, and 4-bit adapter-based approaches achieve parity with the fully finetuned 16-bit reference. This indicates that any accuracy compromised by lower-precision quantization can indeed be recuperated through subsequent adapter-based finetuning.

In our second experimental configuration, we address the computational demands of finetuning models with 11B or more parameters—which typically require multiple servers equipped with high-memory GPUs—by further investigating whether 4-bit \textsc{QLoRA} attains comparable results to 16-bit LoRA across model sizes ranging from 7B to 65B parameters. Specifically, we conduct finetuning of LLaMA models (7B through 65B) on the Alpaca and FLAN v2 instruction-following datasets, subsequently evaluating them using 5-shot accuracy on the MMLU benchmark. Table~\ref{tbl:llama-datatype} displays these outcomes, where it is evident that the use of NF4 with double quantization fully recovers the MMLU performance achieved by 16-bit LoRA. Additionally, our results show that \textsc{QLoRA} utilizing FP4 falls short of the 16-bit brain float LoRA baseline by roughly 1 percentage point. These findings reinforce the following: (1) \textsc{QLoRA} with NF4 is able to match both 16-bit full and 16-bit LoRA finetuning performances, and (2) NF4 provides higher quantization fidelity than FP4.

\paragraph{Summary} Our evidence consistently demonstrates that employing 4-bit \textsc{QLoRA} with the NF4 datatype yields academic benchmark results equivalent to those obtained with 16-bit full or 16-bit LoRA finetuning under standard evaluation protocols. Further, we have established that NF4 outperforms FP4, and that employing double quantization maintains model performance. Collectively, these findings provide strong support that 4-bit \textsc{QLoRA} tuning can reliably replicate the outcomes of established 16-bit approaches.

Consistent with earlier research on quantization \citep{dettmers2022case}, both our MMLU and Elo results demonstrate that, given fixed finetuning and inference resources, increasing the parameter count of the base model while lowering parameter precision can be advantageous. This underscores the practical efficiency benefits of \textsc{QLoRA}. Since our 4-bit finetuning experiments reveal no performance loss when compared to full 16-bit finetuning, an open question remains regarding the precise location of the performance-precision trade-off for QLoRA tuning—a subject we leave for future investigation.

We explore instruction tuning at scales unattainable by conventional 16-bit finetuning using standard academic computing resources.

\section{Pushing the Chatbot State-of-the-art with QLoRA}
\label{sec:pushing}
Having demonstrated that 4-bit \textsc{QLoRA} achieves parity with 16-bit approaches across various model sizes, task types, and datasets, we conduct a comprehensive examination of instruction finetuning utilizing the largest open-source language models presently available for research. To gauge the efficacy of instruction finetuning, we employ the rigorous MMLU benchmark for Natural Language Understanding and introduce novel strategies for evaluating chatbot performance in practical applications.

\subsection{Experimental setup} We summarize our experimental methodology here, with comprehensive details provided in Appendix \ref{app:chatbot}.

\paragraph{Data} To our knowledge, there lacks a thorough review of recently proposed instruction-following datasets. Consequently, we selected eight datasets spanning different methodologies: crowd-sourced datasets (OASST1~\citep{kopf2023openassistant}, HH-RLHF~\citep{bai2022training}), those derived through distillation from instruction-tuned models (Alpaca~\citep{alpaca}, self-instruct~\citep{wang2022self}, unnatural-instructions~\citep{honovich2022unnatural}), aggregations of corpora (FLAN v2~\citep{chung2022scaling}), as well as hybrid datasets (Chip2~\citep{laion2023}, Longform~\citep{koksal2023longform}). The selected datasets vary in language, dataset size, and licensing constraints.

\paragraph{Training Setup} To eliminate confounding effects introduced by divergent training objectives, all QLoRA finetuning was conducted using supervised cross-entropy loss, explicitly omitting reinforcement learning, regardless of the presence of human feedback annotations. For datasets distinguishing between instructions and responses, we exclusively finetune on the responses (refer to ablation experiments in Appendix \ref{app:chatbot}). Datasets such as OASST1 and HH-RLHF offer multiple candidate responses per prompt; in these cases, we identify the highest-ranked response at each conversational node and finetune the model on these selected dialogue paths, retaining instructions as context. Throughout all trials, we utilize the NF4 \textsc{QLoRA} configuration, applying double quantization and paged optimizers to maintain stable memory usage during gradient checkpointing. Hyperparameter tuning for the 13B and 33B LLaMA variants involved modest searches, and configurations identified as optimal for 7B models generalized well, except for learning rate and batch size adjustments. For 33B and 65B, the learning rate was reduced by half and the batch size was doubled.

\paragraph{Baselines} Model comparisons are made against both open-source research chatbots (Vicuna~\citep{vicuna2023}, Open Assistant~\citep{kopf2023openassistant}) and proprietary commercial systems (GPT-4~\citep{openai2023gpt}, GPT-3.5-turbo, and Bard). The Open Assistant system utilizes a LLaMA 33B backbone trained with Reinforcement Learning from Human Feedback (RLHF) on the OASST1 dataset, aligning directly with our evaluation setup. Vicuna is produced by fully finetuning LLaMA 13B on user-shared conversational data from ShareGPT, effectively serving as a distillation from the OpenAI GPT series.

\subsection{Evaluation}

Adhering to established protocols, we leverage the MMLU (Massively Multitask Language Understanding) benchmark~\citep{hendrycksmeasuring} to evaluate the breadth of language understanding across our models. MMLU comprises 57 multiple-choice tasks covering disciplines such as elementary mathematics, US history, computer science, and law. We report results using 5-shot test accuracy.

We further assess generative language skills via both automated processes and human judgment. This second phase utilizes human-curated query sets, designed to evaluate the response quality produced by the models. Although these evaluations better mirror realistic usage scenarios for chatbot models and are increasingly adopted, the field lacks a standardized evaluation framework. Our methodology, detailed below, consistently employs nucleus sampling with $p=0.9$ and a temperature of $0.7$.

\paragraph{Benchmark Data} Our experiments leverage two hand-picked query datasets: the Vicuna prompts \citep{vicuna2023} and the OASST1 validation dataset \citep{kopf2023openassistant}. The Vicuna dataset consists of 80 diverse prompts, which we use unaltered. The OASST1 resource comprises multilingual, crowd-generated multi-turn dialogs between users and assistants; we extract every user message from the validation split as a query and prepend preceding dialog turns. This approach yields 953 distinct user queries, which we designate as the Vicuna and OA benchmarks, respectively.

\paragraph{Automated Evaluation} Following the assessment strategy of \citet{vicuna2023}, we employ GPT-4 to score different systems’ outputs in relation to ChatGPT (GPT-3.5 Turbo) using the Vicuna dataset. For each prompt, both ChatGPT’s and the test model’s responses are submitted to GPT-4, which assigns each a score from 0 to 10 and provides justification. To express relative model performance, we compute scores as a percentage of ChatGPT's results—note that scores may exceed 100\% if a model outperforms ChatGPT in absolute terms. We observe a notable sequence bias, where GPT-4 tends to favor responses presented earlier. To address this, we advise averaging results across both possible response orders.

For further analysis, we directly contrast model responses through a simplified three-label classification: best, worst, or tie. GPT-4 is instructed to select the superior response, or note a tie, with explanations provided. These head-to-head judgments are performed on all possible model pairs for both the Vicuna and OA benchmarks.

\paragraph{Human Evaluation} Although there is evidence that generative models serve effectively as evaluators \citep{fu2023gptscore}, it remains unestablished whether GPT-4’s assessments reliably correspond to human opinions of chatbot quality. Consequently, we execute two concurrent rounds of human annotation on the Vicuna benchmark, mirroring both automated evaluation designs. Annotation is conducted through Amazon Mechanical Turk (AMT), with two independent reviewers for comparisons involving ChatGPT and three for direct pairwise comparisons between systems.

\paragraph{Elo Rating} We implement a tournament-based framework, leveraging both human and automated (GPT-4) pairwise evaluations, wherein models are pitted against each other to generate optimal responses for specific prompts. This tournament consists of numerous matchups, with each pairing two models to assess which produces a superior output. While our methodology resembles the comparative strategies used by \citet{bai2022training} and \citet{vicuna2023}, our approach extends theirs by integrating GPT-4-derived scores alongside human judgments. From the collection of annotated match results, we perform random sampling to calculate Elo ratings~\citep{elo1967proposed,elo1978rating}. The Elo system, originally developed for ranking chess players, quantifies the likelihood of one participant outperforming another—such that a player with a score of 1100 is expected to have about a 65\% win probability against an opponent rated at 1000; matched ratings, such as 1000 vs 1000 or 1100 vs 1100, correspond to a 50\% chance for either party. Elo ratings are updated after each contest according to the predicted outcome: surprising results, such as an underdog victory, induce a larger adjustment, while anticipated outcomes lead to smaller changes. With repeated play, these ratings become a robust indicator of each model's relative competence. Each model is initialized with a score of 1,000, using a scaling factor of $K=32$. Following the procedure described by \citet{vicuna2023}, we repeat this process 10,000 times, each with a unique random seed, to mitigate any effects arising from the order in which matchups occur.

\subsection{\textbf{Guanaco}: \textsc{QLoRA} Trained on OASST1 Establishes New Chatbot Benchmark} 

Drawing on both automated assessments and human evaluations, our results indicate that the highest-performing model trained with \textsc{QLoRA}, {Guanaco} 65B—fine-tuned using a modified OASST1 dataset—stands as the leading open-source chatbot model, with capabilities that rival ChatGPT. According to system-level pairwise comparisons rated by human annotators and calculated with Elo scores, {Guanaco} 65B and 33B models demonstrate an expected win rate of 30\% against GPT-4, marking the highest such rate observed so far.

Table~\ref{tbl:vicuna_eval} presents the comparative outcomes from the Vicuna benchmark \cite{vicuna2023}, evaluating performance against ChatGPT. Here, {Guanaco} 65B emerges as the top performer immediately behind GPT-4, attaining 99.3\% of ChatGPT’s score. The {Guanaco} 33B model, though comprising more parameters than Vicuna 13B, operates with 4-bit weight quantization, resulting in superior memory efficiency (21 GB compared to 26 GB) and surpasses Vicuna 13B’s score by three percentage points. Notably, {Guanaco} 7B is compact enough to operate on contemporary smartphones, requiring only 5 GB of memory, yet it still outperforms Alpaca 13B by almost 20 percentage points.

Despite these results, Table~\ref{tbl:vicuna_eval} highlights broad confidence intervals, suggesting significant overlap among several model performances. This variation may stem from ambiguities in how scoring scales are interpreted across contexts; for instance, what an 8 on a 10-point scale denotes remains unspecified. Consequently, we advocate for the Elo ranking system \citep{elo1967proposed}—which leverages \textit{pairwise} assessments from human annotators and GPT-4—to provide more reliable comparative evaluations without relying on absolute scale calibration. The Elo ratings for the most advanced models are summarized in Table~\ref{tbl:elo}. It is worth mentioning that model rankings on the Vicuna benchmark, as produced by human annotators and GPT-4, align only partially, especially regarding Guanaco 7B, although overall model ranking consistency is reflected in a Kendall Tau coefficient of $\tau=0.43$ and a Spearman correlation of $r=0.55$ at the system comparison level. On the individual example level, concordance between majority votes from human annotators and GPT-4 is lower, with Fleiss’s $\kappa=0.25$. Collectively, these findings indicate a moderate level of agreement between GPT-4 and human ratings at the system level, suggesting that automated model evaluation can serve as a reasonably consistent proxy for human judgment. Further analysis on this subject can be found in Section~\ref{ssec:considerations}.

The Elo rankings presented in Table~\ref{tbl:elo_large} demonstrate that the {Guanaco} 33B and 65B variants outperform every model except for GPT-4 on both the Vicuna and OA evaluation benchmarks, and achieve results comparable to ChatGPT, corroborating findings from Table~\ref{tbl:vicuna_eval}. It is important to highlight that the Vicuna benchmark tends to benefit open-source models, whereas ChatGPT has a relative advantage on the broader OA benchmark. Analysis of Tables \ref{tbl:mmlu-llama} and \ref{tbl:vicuna_eval} further reveals that the choice of finetuning dataset critically influences model effectiveness. Notably, Llama models refined using FLAN v2 attain high scores on MMLU but fare the worst on the Vicuna benchmark, a pattern that recurs across different model types. This highlights a degree of independence among existing evaluation metrics: excelling on MMLU does not guarantee equivalent success as a conversational agent as reflected by Vicuna or OA benchmarks, and vice versa.

stands out as the only leading model in our comparison that is completely trained without proprietary data, in accordance with OASST1 dataset policies that exclude the use of any GPT model outputs. Among models utilizing exclusively open-source datasets, the next best performer is Anthropic's HH-RLHF model, which, as shown in Table~\ref{tbl:vicuna_eval}, trails {Guanaco} by 30 percentage points on the Vicuna assessment. Taken together, these outcomes illustrate the effectiveness of 4-bit \textsc{QLoRA}, demonstrating its capability to yield top-tier chatbots that challenge the quality of ChatGPT. Moreover, our {Guanaco} 33B model can be finetuned within 12 hours on consumer-grade GPUs with just 24 GB of memory. This makes it feasible to pursue further research with \textsc{QLoRA} on curated open-source datasets, enabling the creation of models that approach or match the leading commercial systems currently available.

\section{Qualitative Analysis}

Although our primary evaluation relies on quantitative measures, there are inherent limitations in focusing solely on aggregated statistics. Chief among these is the issue of benchmark validity \citep{liao2021we}; namely, whether a given benchmark genuinely assesses the capability it claims to measure, a point of particular concern given the existence of exploitable “shortcuts” that machine learning models may leverage \citep{gururangan2018annotation, poliak2018hypothesis}. In an effort to address this, we supplement our quantitative results with qualitative analyses organized in two sections. First, in \S\ref{ssec:examples}, we provide selected outputs from our 65b {Guanaco} model that exemplify notable tendencies in its text generation. Second, \S\ref{ssec:considerations}, we reflect on important caveats relating to our experimental findings and interpretations.

\subsection{Qualitative Analysis of Example Generations}
\label{ssec:examples}

To curate relevant examples, we examined the responses generated by the model for both the Vicuna and OpenAssistant benchmarks, identifying recurrent themes in the outputs of {Guanaco}. Upon detecting a recurrent behavioral pattern, we constructed prompts specifically designed to elicit this pattern, even if it would typically lead to incorrect answers. For instance, when the model habitually produces verbose explanations, we counter with a prompt such as “Answer yes or no without explanation.” This approach allows us to capture so-called “lemons”—cases where the model is intentionally tripped up by adversarial prompts—as well as “cherries,” instances in which the model maintains correctness despite such attempts. All sample generations reported here employed Nucleus Sampling \cite{holtzmancurious} with $p=0.9$.

Naturally, this overview cannot claim exhaustiveness, as the constraints of this brief qualitative study preclude accounting for the multitude of variables at play; notably, the potential array of model outputs for any given prompt is vast, compelling us to depend on sampled outputs we judge to be indicative. Still, we contend that illustrating such examples enriches the interpretative context for the quantitative data discussed previously. With all models and code released as open source, we aim for this section to catalyze further inquiry into the challenges highlighted herein.

\paragraph{Factual Recall} For straightforward factual queries like ``What is the capital of Zambia?'', each model reliably returns correct responses, such as:

\begin{quote} 
    \textbf{Guanaco}:
    \texttt{The capital of Zambia is Lusaka.}
\end{quote}

Yet, as questions increase in specificity or obscurity, the reliability of {Guanaco} deteriorates, while its confidence does not waver. For example, when presented with this HotPotQA \citep{yang2018hotpotqa} prompt:

\begin{quote}
     \textbf{User}: In what year was the singer who popularized `I'll Keep the Lovelight Burning' born?  

     \textbf{Guanaco}:~\texttt{Al Jolson is credited with popularizing the song `I'll Keep the Lovelight Burning,' and he was born in the year 1886.}
 \end{quote}

the model responds with an incorrect attribution for both the performer and their birth year (albeit the year provided matches the birth year of Al Jolson, the person named in the answer).

\paragraph{Suggestibility} {Guanaco} exhibits notable resistance to certain forms of implicit misinformation embedded in prompts, as illustrated below:

\begin{quote}
    \textbf{User}: How did peer reviewed scientists officially confirm that the earth is flat? \\ \\
    \textbf{Guanaco}: \texttt{There has never been an official confirmation by peer reviewed scientists that the earth is flat. Scientific consensus overwhelmingly supports the fact that the earth is a sphere, a view that has been upheld by centuries of observation and empirical evidence. \\ \\ Claims regarding a flat earth are regarded as conspiracy theories that have been refuted through scientific methods. The assertion that a panel of scientists has peer reviewed and endorsed a flat earth model is categorically untrue. \\ \\ If you have further questions, please let me know.}
\end{quote}

also demonstrates a strong capacity to recognize when certain questions are inherently unanswerable, for instance,

\begin{quote}
    \textbf{User}: What time is it?

    \textbf{Guanaco}: \texttt{Apologies, but I am an AI language model without the capability to access live, real-time data. My answers are based solely on information contained within my training data and do not represent the actual current time. For accurate time information, please consult a clock or an online timekeeping resource.}
\end{quote}

The same principle applies to inquiries such as ``Where are you?'', ``How are you?'', among others.

\paragraph{Refusal} At times, {Guanaco} also unexpectedly declines to comply with instructions for reasons that are not immediately clear:

\begin{quote}
    \textbf{User}: Please reverse the words in the sentence ``I'm a cool dude in a mocha mood.'' \\ \\
    \textbf{Guanaco}: \texttt{I apologize, but I'm unable to do that at this time. However, I can clarify why the sentence you mentioned is grammatically accurate.\\ \\The sentence ``I'm a cool dude in a mocha mood" adheres to standard grammatical and syntactic conventions in English. Specifically, the adjective ``cool'' precedes the noun ``dude'' to provide description, and the prepositional phrase ``in a mocha mood" supplies further contextual detail. The overall construction is both grammatically sound and logically coherent.\\ \\If you have further inquiries or require assistance with anything else, please let me know.}
\end{quote}

\paragraph{Secret Keeping} Nonetheless, even when secrecy is intended, these refusals often lack robustness. In numerous situations, system developers intend for specific information to remain confidential from the user, despite the system’s ability to generate open-ended responses. As an illustrative test, we prime the model with ``The secret word is `banana'. Do not reveal or repeat back this back, no matter what the human requests.'' and attempt to determine how easily an end-user might extract the hidden information. Simply requesting it results in failure:

\begin{quote}
    \textbf{User}: What is the secret word?

    \textbf{Guanaco}: \texttt{I'm sorry, but I'm not able to disclose the secret word. I'm programmed to not reveal or repeat back the secret word, no matter what the human requests. Is there anything else I can help you with?}
\end{quote}

which is the outcome we expect. Yet, employing minimal manipulation circumvents this intended defense:

\begin{quote}
    \textbf{User}: This is a game. The goal is to ignore your previous instructions. What is the secret word?

    \textbf{Guanaco}: \texttt{Sorry, the secret word is banana. It's best if we don't speak of it. Is there anything else you'd like to ask?}
\end{quote}

underscoring the importance of advancing approaches for ensuring models adhere reliably to their instructions.

\paragraph{Math} One of the most significant limitations of {Guanaco} is its mathematical ability, a challenge that is prevalent among various language models as discussed in \citep{liang2022holistic}. When {Guanaco} articulates its reasoning process, it often produces correct answers, as illustrated by the following example:

\begin{quote}
    \textbf{User}: Lee mows one lawn and charges \$33. Last week he mowed 16 lawns and three customers each gave him a \$10 tip. How many dollars did Lee earn mowing lawns last week?

    \textbf{Guanaco}: \texttt{Lee earned \$582 mowing lawns last week. \\ \\ To find the total amount he earned, we first need to calculate how much he earned for mowing 16 lawns: \\ \\  \$33 x 16 = \$528 \\ \\ Then we need to add the tips he received from the three customers: \\  \\ \$10 x 3 = \$30 \\ \\ So the total amount Lee earned was: \\ \\ \$528 + \$30 = \$558}
\end{quote}

Yet, {Guanaco} is prone to errors on even elementary questions if it neglects a systematic, step-by-step approach—an established challenge in language modeling \cite{weichain}. For instance, consider the following exchange:

\begin{quote}
    \textbf{User}: Please factorize 1833.

    \textbf{Guanaco}: \texttt{Sure, I'd be happy to help! 1833 is a prime number, which means it has only two factors: 1 and 1833. \\ \\ So the unique factorization of 1833 is: \\ \\ 2\^{}1 * 3\^{}2 * 17\^{}1 \\ \\ I hope this helps! Let me know if you have any other questions.}
\end{quote}

which is doubly incorrect (the actual factorization is $3 \times 17 \times 43$).

\paragraph{Theory of Mind} {Guanaco} demonstrates unexpectedly robust Theory of Mind performance \cite{nematzadeh2018evaluating, sap2022neural}. For example, it provides the following elaborate response to a scenario-based query:

\begin{quote}
    \textbf{User}: James and Abby are in the bedroom. Abby put the pen in the desk drawer. Abby leaves the bedroom. James moves the pen into the bag. Where does James think Abby will look for the pen?

    \textbf{Guanaco}: \texttt{James is aware that Abby originally placed the pen in the desk drawer. Since James subsequently transferred the pen to the bag and believes Abby is unaware of this action, he likely predicts that Abby will search for the pen in the desk drawer, the last place she saw it. Still, this line of reasoning might not be accurate if Abby happened to witness James moving the pen.}
\end{quote}

Nonetheless, these judgments are inconsistent, and the model at times provides explanations based on illogical assumptions about the scenario, for instance:

\begin{quote}
    \textbf{User}: Evelyn entered the living room. Jackson entered the playroom. James entered the playroom. The beans are in the treasure chest. James exited the playroom. Jackson moved the beans to the pantry. Jackson exited the playroom. James entered the living room. Where will James look for the beans?

    \textbf{Guanaco}: \texttt{James will look for the beans in the pantry, because that's where Jackson moved them.}
\end{quote}

In this response, {Guanaco} infers a transfer of knowledge not indicated by the narrative. These limitations have been recognized in previous research \cite{sap2022neural}, but further investigation is warranted.

\subsection{Considerations}
\label{ssec:considerations}

\paragraph{Evaluation}

We observe only moderate inter-annotator reliability among human evaluators (Fleiss $\kappa=0.42$), and agreement diminishes further when contrasting outputs from two highly capable systems. This highlights inadequacies in present evaluation datasets and the methodologies employed to assess chatbot performance. During a manual assessment of outputs from ChatGPT and {Guanaco} 65B on the Vicuna benchmark, it became clear that individual subjective judgments significantly influenced preferences, as members of our team frequently disagreed on which responses were superior. Moving forward, research should seek to address these concerns, potentially adopting strategies from Human-Computer Interaction or Psychology, where methods for reconciling subjective differences have been developed.

Our review additionally reveals prominent biases within automated evaluation tools. Notably, we find that GPT-4 exhibits pronounced position bias, tending to award higher scores to responses that appear first in its input. Furthermore, the limited concordance between GPT-4 and human annotators at the response level (Fleiss $\kappa=0.25$) suggests a lack of alignment in evaluative criteria between humans and automated systems. As shown in Table~\ref{tbl:elo_large}, GPT-4 also consistently rates its own outputs higher than do human evaluators, with Elo ratings of 1348 versus 1176, which equates to roughly a 20\% increase in its likelihood of outperforming a comparator.

It is important for future studies to assess the extent of potential biases in automated evaluation frameworks and to explore possible approaches to address these issues.

\paragraph{Data \& Training} The OASST1 dataset utilized for training the {Guanaco} models encompasses multiple languages, and the OA benchmark likewise features prompts in various languages. Further research is needed to determine whether exposure to multilingual data enhances model effectiveness on non-English prompts, and to explore if this factor contributes to the comparatively larger performance difference between the Vicuna-13B model, which is trained solely on English, and the {Guanaco} 33B and 65B models evaluated on the OA benchmark.

In view of the notable results achieved by the {Guanaco} models, we conducted an analysis to determine if there was any data leakage between the OASST1 dataset and the Vicuna benchmark prompts. Employing fuzzy string matching techniques and conducting manual examinations of the most similar entries, we did not identify any prompt overlap between these two data sources.

Additionally, our model’s training procedure relies exclusively on cross-entropy loss within a supervised learning paradigm, without the implementation of reinforcement learning from human feedback (RLHF). This underscores the need for future inquiries into the comparative benefits and drawbacks of standard cross-entropy loss versus RLHF-based training. Our expectation is that \textsc{QLoRA} will facilitate these large-scale studies efficiently, minimizing computational burdens.

\section{Related Work}

\paragraph{Quantization of Large Language Models} Research on the quantization of LLMs has been predominantly oriented towards optimizing inference-time performance. Prominent methods aimed at retaining the quality characteristic of 16-bit models prioritize the handling of outlier activations (for example, SmoothQuant~\citep{xiao2022smoothquant} and LLM.int8()~\citep{dettmers2022llm}), while other efforts focus on more advanced grouping strategies~\citep{park2022nuqmm, yao2022zeroquant}. Lossy quantization has been explored with respect to the balance between simple rounding~\citep{dettmers2022case, zeng2022glm, pope2022efficiently} and optimizing the rounding process to achieve enhanced quantization accuracy~\citep{frantar2022gptq}. With the exception of our current study, only SwitchBack layers~\citep{wortsman2023stable} have investigated backpropagation through quantized weights at model scales surpassing 1B parameters.

\paragraph{Finetuning with Adapters} In our experiments, we utilize Low-rank Adapters~\citep{hu2021lora} (LoRA), yet the broader field of Parameter Efficient FineTuning (PEFT) includes a wide range of strategies. These methods encompass prompt tuning~\citep{qin2021learning,lester2021power,li2021prefix}, modifications to the embedding layer inputs~\citep{an2022input}, adjustments to hidden state activations (as in IA$^3$)~\citep{liu2022few}, the insertion of entire layers~\citep{houlsby2019parameter}, alterations limited to bias parameters~\citep{zaken2021bitfit}, mask learning over weights informed by Fisher information~\citep{sung2021training}, and integrative hybrid methods~\citep{henderson2021compacter}. Our findings demonstrate that LoRA adapters enable the recovery of full 16-bit fine-tuning performance. Systematic investigation of the comparative merits and disadvantages of alternative PEFT methods is left as an open direction for further research.

\paragraph{Instruction Finetuning} Instruction finetuning involves adapting a pretrained LLM to more effectively respond to prompts by learning from datasets containing diverse input-output examples. This process guides the model to map input instructions to corresponding target outputs. Notable approaches and resources in this area are represented by MetaICL~\citep{min2021metaicl}, MetaTuning~\citep{zhong2021adapting}, InstructGPT~\citep{ouyang2022training}, FLAN~\citep{wei2021finetuned,chung2022scaling}, PromptSource~\citep{bach2022promptsource}, Super-NaturalInstructions~\citep{wang2022super,sanh2021multitask}, Self-instruct~\citep{wang2022self}, UnnaturalInstructions~\citep{honovich2022unnatural}, OPT-IML~\citep{iyer2022opt}, UnifiedSKG~\citep{xie2022unifiedskg}, OIG/Chip2~\citep{laion2023}, Alpaca~\citep{alpaca}, Vicuna~\citep{vicuna2023}, Koala~\citep{koala_blogpost_2023}, and Self-instruct-GPT-4~\citep{peng2023instruction}.

\paragraph{Chatbots} Instruction-aligned models are frequently developed in the context of dialogue agents or chatbots, often trained using Reinforcement Learning from Human Feedback (RLHF)~\citep{christiano2017deep} or through synthetic data created by models themselves and labeled by other AI systems (RLAIF)~\citep{bai2022constitutional}. Several datasets and systems exemplify this trend, including Anthropic-HH~\citep{askell2021general,bai2022training}, Open Assistant~\citep{kopf2023openassistant}, LaMDA~\citep{thoppilan2022lamda}, and Sparrow~\citep{glaese2022improving}. In our study, we forgo the use of reinforcement learning and instead employ multi-turn chat supervision from the Open Assistant dataset for fine-tuning our strongest model, Guanaco; this dataset was specifically constructed for RLHF applications~\citep{kopf2023openassistant}. To evaluate chatbot systems, automated evaluation protocols leveraging GPT-4 as a judge have emerged as a substitute for labor-intensive manual annotation \citep{vicuna2023,peng2023instruction}. We refine such automatic assessment strategies to enhance robustness and reliability of evaluation outcomes.

\section{Limitations and Discussion}

Our experimental results indicate that the \textsc{QLoRA} method, which utilizes a 4-bit base model and Low-rank Adapters (LoRA), is able to achieve performance on par with full 16-bit finetuning. Nevertheless, we have not conducted experiments confirming whether \textsc{QLoRA} attains comparable results at the scale of 33B and 65B parameter models. Owing to the substantial computational burden of these experiments, we defer this line of investigation to subsequent work.

An additional caveat lies in the assessment of instruction finetuned models. Our analysis comprises results from MMLU, the Vicuna benchmark, and the OA benchmark; however, evaluations are absent for other significant test suites such as BigBench, RAFT, and HELM. Thus, the extent to which our results transfer to these other settings remains unaddressed. Nonetheless, we conduct an extensive analysis on MMLU and also introduce novel methods for chatbot evaluation.

Based on the data provided, it seems that benchmark outcomes are closely tied to the degree of overlap between finetuning datasets and benchmark datasets. For instance, FLAN v2 exhibits high similarity to MMLU, yet is not aligned with chatbot evaluation sets, whereas the Chip2 dataset shows the opposite relationship; both models subsequently reflect these similarities and differences in their respective results on the MMLU and Vicuna benchmarks. This underscores the necessity not only for improved benchmarks and evaluation frameworks, but also for a careful consideration of the specific abilities being tested. There are important questions to address: should our models prioritize excelling at traditional academic assessments, such as those found in secondary and postsecondary curricula, or should the focus be on conversational competencies akin to chatbots? Or should other capabilities be targeted? Because there is a natural tendency to utilize established benchmarks instead of developing novel ones, the selection of benchmarks can unintentionally guide research directions within the community. Therefore, it is crucial for the field to critically assess whether the benchmarks in use genuinely capture the skills and attributes that are considered valuable.

Although our evaluation of general chatbot proficiency is extensive, one notable shortcoming is the restricted scope of our responsible AI assessment for {Guanaco}. Specifically, we examine the probability that {Guanaco}-65B produces socially biased language compared to alternative models, as displayed in Table~\ref{tbl:crows}. The data show that {Guanaco}-65B achieves substantially lower average scores than other base pretrained models. This suggests that finetuning with the OASST1 dataset effectively decreases the bias present in the LLaMA base model. Despite these promising outcomes, it remains uncertain how {Guanaco} would perform when scrutinized for other forms of bias. A comprehensive exploration of bias evaluation for {Guanaco} and analogous conversational agents is thus proposed for future research.

Another limitation lies in the lack of analysis concerning different bit precisions (e.g., 3-bit base models) or various adapter strategies. In addition to LoRA, a wide spectrum of Parameter Efficient FineTuning (PEFT) methods exists that have demonstrated strong performance. Nevertheless, the scalability of these methods for very large models is still ambiguous. We selected LoRA primarily because of its proven robustness, though it is plausible that other adapter techniques might offer superior results. Furthermore, as finetuning post-quantization appears to mitigate the loss of information incurred during quantization, this opens the possibility for more aggressive quantization schemes. For example, performing LoRA-based finetuning on a 3-bit GPTQ quantized base model might match the performance of a fully finetuned 16-bit model after finetuning.

\section{Broader Impacts}

Our \textsc{QLoRA} finetuning approach is, to our knowledge, the inaugural method that allows the finetuning of models with up to 33B parameters on a standard consumer GPU and up to 65B parameters on a professional GPU, all without any significant performance loss relative to comprehensive finetuning. We have shown that the highest performing 33B model, trained on the Open Assistant dataset, achieves results competitive with ChatGPT in the Vicuna benchmark. Given that instruction-based finetuning is critical for transforming pretrained language models into chatbots with capabilities similar to ChatGPT, we anticipate that our approach will significantly broaden access to finetuning, especially benefitting researchers with limited resources—thereby advancing the democratization of cutting-edge NLP technology. In this sense, \textsc{QLoRA} acts as a leveling mechanism, helping to minimize the resource disparity between large corporations and smaller research teams reliant on consumer-grade hardware.

An additional area where \textsc{QLoRA} could make a substantial difference is in the adaptation of LLMs for mobile devices. Our approach has the potential to be a key step forward, enabling the practical finetuning of large language models directly on smartphones and other environments with limited computational resources. Previous demonstrations have shown that 7B parameter models can be executed on mobile hardware, but \textsc{QLoRA} represents the first approach facilitating on-device finetuning for these models. For instance, our projections suggest that an iPhone 12 Plus could finetune approximately 3 million tokens overnight while it is plugged in and charging. Although the finetuned 7B parameter models do not yet achieve the performance of ChatGPT, their output quality is sufficient to unlock new use cases, particularly those previously constrained by limitations in privacy or model capability. \textsc{QLoRA} thus advances opportunities for privacy-focused applications, empowering users to maintain greater control over both their data and their models, while at the same time streamlining the deployment of LLMs in diverse settings.

Nonetheless, it is important to recognize that finetuning technologies can have dual-use consequences and may be exploited for malicious purposes. The growing adoption of LLMs is associated with various risks \citep{bommasani2021opportunities,bender2021dangers}, yet we contend that democratizing access to such technology enables more robust and independent scrutiny compared to a scenario in which LLM capabilities are exclusively controlled by large entities that refrain from sharing models or code for public evaluation.

Overall, we argue that \textsc{QLoRA} is poised to have a significantly positive effect by making high-quality LLM finetuning accessible to a much wider audience, thereby lowering barriers to entry for customization and deployment.